\subsection{Interpretation of the Hierarchical Model under Strict Protocols}

The hierarchy separates two engineering questions that have different label spaces: whether a telemetry window is abnormal and which labeled fault domain best describes an abnormal window. This separation makes the decision path explicit and permits fault-domain-specific evidence to be generated only after an anomaly decision. It also clarifies why Stage 2 performance and complete Cascade performance differ. Stage 2 is evaluated on anomaly-only windows, whereas the Cascade inherits false negatives and false positives from the detector before applying the fault-domain classifier.

The protocol comparisons show that this interpretability advantage should not be conflated with predictive superiority. On the fixed chronological split, the paired Direct-minus-Cascade Macro-F1 contrast was $-0.0303$ (95\% flight-log-bootstrap CI $[-0.0524,-0.0086]$), favoring the Cascade for that split. Purging reversed the mean ordering, but its paired interval included zero; the leave-log-out interval likewise included zero. The available evidence thus does not support a protocol-independent advantage for either complete five-class formulation. The contribution of the hierarchy lies instead in its explicit gating semantics, its separation of detector and diagnostic errors, and its traceable explanation structure.

The degradation from the fixed chronological protocol to stricter isolation also has methodological implications. Overlapping windows from the same flight can share local operating conditions even when assigned chronologically, and a boundary purge removes only part of this proximity. Assigning entire flight logs to disjoint partitions provides a stronger estimate of transfer to unseen flights. The lower leave-log-out Stage 1 AUPRC and event F1 indicate that the detector is the principal bottleneck in the complete chain, whereas the stable anomaly-only Stage 2 result suggests that labeled fault-domain structure remains learnable once abnormal windows are supplied. Future model comparisons should therefore retain flight-log-level isolation and report detector, conditional diagnostic, and complete-task metrics separately.

The full-channel CATCH execution closes the earlier input-mismatch concern but not every comparability limitation. It used the official model core, the same 87 raw channels and chronological windows, three seeds, and both window- and event-level metrics; however, its normal-only objective, 4096-window training sample, and two-epoch budget differ intrinsically from the supervised LightGBM protocol. The GCAD-inspired result additionally remains a non-official linear ablation. These references therefore show that the evaluated recent detectors were weak under the declared telemetry representation and budgets, not that HS-AeroTS is unconditionally superior to the published methods.
